% ====================================================================
% === SUPERSEDED (reader-facing) by Docs/notes/12_spot_vol_dynamics.tex
%     (Vol-Fitter Technical Notes series). For the maintained,
%     comprehensive treatment with implementation + performance
%     appendices and code-generated figures, read that note.
% === RETAINED here as the equation-label citation source for code
%     docstrings, which reference this file by path and by its eq.
%     labels (e.g. "eq. q_logit", "note section 5.2").
% ====================================================================
\documentclass[8pt]{extarticle}
\usepackage[a4paper,margin=0.44in]{geometry}
\usepackage{amsmath,amssymb,mathtools,booktabs,array}
\usepackage{enumitem}
\usepackage{microtype}
\usepackage[compact]{titlesec}
\titlespacing*{\section}{0pt}{4pt}{1pt}
\setlength{\parindent}{0pt}
\setlength{\parskip}{1.4pt}
\setlist[itemize]{leftmargin=1.1em,itemsep=0.5pt,topsep=1pt}
\setlist[enumerate]{leftmargin=1.3em,itemsep=0.5pt,topsep=1pt}
\newcommand{\sig}{\sigma}
\newcommand{\w}{w}
\newcommand{\ATM}{\mathrm{ATM}}
\newcommand{\loc}{\mathrm{loc}}
\newcommand{\str}{\mathrm{str}}
\newcommand{\mon}{\mathrm{mon}}
\newcommand{\LV}{\mathrm{LV}}
\newcommand{\e}{\mathrm{e}}
\newcommand{\eps}{\varepsilon}
\begin{document}
\begin{center}
{\large\bfseries Fast Spot-Move Updates of an Implied-Volatility Surface}\vspace{1pt}

{\small Sticky-strike, sticky-moneyness, Hagan sticky-local-vol, and SSR, including local-vol-grid rules}\vspace{1pt}
\end{center}

\textbf{Common coordinates.}
For maturity $T$, let $F_T^0$ be the old forward and $F_T^1$ the new forward after the spot move. Define
\[
        h_T:=\log\frac{F_T^1}{F_T^0},\qquad
        \Delta F_T:=F_T^1-F_T^0,\qquad
        k:=\log\frac{K}{F_T^1} .
\]
Thus $k$ is the \emph{new} log-forward moneyness, while the same strike $K=F_T^1\e^k$ had old moneyness
\[
        k^0(k)=\log\frac{K}{F_T^0}=k+h_T .
\]
Let
\[
        \sig_0(T,k),\qquad \w_0(T,k):=T\sig_0^2(T,k),\qquad
        \kappa_T:=\partial_k\sig_0(T,0),\qquad
        g_T:=\partial_k\w_0(T,0)=2T\sig_0(T,0)\kappa_T .
\]
In production, transport/interpolate total variance $\w$ and then set $\sig=\sqrt{\w/T}$. For the local-vol grid, denote the old grid/interpolator by
\[
        L_0(t,K):=\sig_{\loc,0}(t,K),
        \qquad h_t:=\log(F_t^1/F_t^0),\qquad \Delta F_t:=F_t^1-F_t^0.
\]
If rates, dividends, borrow, or discrete dividends are affected, all $h_t,h_T$ must be computed from forwards, not directly from the spot ratio.

\section*{1. Sticky-strike}
Sticky-strike means that the implied volatility attached to an absolute strike is unchanged:
\[
        \sig_1^{\str}(T,K)=\sig_0(T,K).
\]
In the new normalized coordinate $k$, the full-smile update is therefore
\[
        \boxed{\;\w_1^{\str}(T,k)=\w_0(T,k+h_T)\;},
        \qquad
        \sig_1^{\str}(T,k)=\sqrt{\w_0(T,k+h_T)/T} .
\]
Thus an upward spot/forward move, $h_T>0$, translates the old smile left when plotted against new log-moneyness. At the new ATM point,
\[
        \sig_{\ATM,1}^{\str}(T)=\sig_0(T,h_T)
        =\sig_0(T,0)+\kappa_T h_T+O(h_T^2).
\]
For an equity-style negative skew, $\kappa_T<0$, spot up gives ATM volatility down; for a fixed strike $K$, volatility is unchanged by construction.

\textit{Local-vol grid.} The exact Dupire local-vol grid consistent with a sticky-strike \emph{implied} surface is not obtained by merely keeping $L_0$ fixed; it would require recomputing local vol from the shifted option-price surface under the new forward. A fast no-recalibration proxy is the \emph{half-grid transport}
\[
        \boxed{\;L_1^{\str}(t,K)\approx L_0\bigl(t,K\e^{-h_t/2}\bigr)\;},
        \qquad K_i^{1,\str}(t)=K_i^0(t)\e^{h_t/2} .
\]
It is the geometric midpoint between the fixed absolute-strike grid and the grid moving fully with the forward. The motivation is the Hagan/weak-skew rule: a local-vol grid shifted by a fraction $\beta$ of the forward move gives an ATM stickiness ratio $R\approx2(1-\beta)$; sticky-strike has $R=1$, hence $\beta=1/2$. Near ATM, the equivalent additive Hagan form is $K_i^{1,\str}\approx K_i^0+\tfrac12\Delta F_t$. Use the exact sticky-strike formula for vanilla IVs, and this half-grid rule only when a refreshed LV grid is needed cheaply.

\section*{2. Sticky-moneyness / sticky-delta}
Sticky-moneyness means that the smile is invariant as a function of log-forward moneyness:
\[
        \boxed{\;\w_1^{\mon}(T,k)=\w_0(T,k)\;},
        \qquad
        \sig_1^{\mon}(T,k)=\sig_0(T,k) .
\]
Consequently,
\[
        \sig_{\ATM,1}^{\mon}(T)=\sig_0(T,0),
        \qquad
        \Delta\sig_{\ATM}^{\mon}=0 .
\]
For a fixed strike with old moneyness $q=\log(K/F_T^0)$, the new moneyness is $q-h_T$, hence
\[
        \sig_1^{\mon}(T,K)=\sig_0(T,q-h_T).
\]
Sticky-delta is the same idea when the smile is parameterized by delta; the delta convention, forward, discounting, and premium adjustment must be consistent.

\textit{Local-vol grid.} The LV grid is unchanged in normalized strike space $x=\log(K/F_t)$. Equivalently, in absolute strike coordinates,
\[
        \boxed{\;L_1^{\mon}(t,K)=L_0\bigl(t,K\e^{-h_t}\bigr)\;},
        \qquad K_i^{1,\mon}(t)=K_i^0(t)\e^{h_t} .
\]
Thus every LV node is scaled by the same forward ratio at its local time. A grid stored directly in $x$ is simply reused; a grid stored in $K$ is relabelled by $K\mapsto K F_t^1/F_t^0$.

\section*{3. Sticky-local-vol / sticky-tree: Hagan local-vol expansion}
The exact sticky-local-vol rule is not an implied-volatility rule. It says: keep the calibrated local-volatility grid fixed in absolute strike/spot level,
\[
        \boxed{\;L_1^{\LV}(t,K)=L_0(t,K)\;},
\]
restart the pricing model from the new spot/forward, and re-invert vanilla prices to Black implied volatility. This is model-faithful, but it is not the fastest update if one only wants a refreshed vanilla surface.

\textit{Local-vol grid.} This is the only one of the four rules for which the LV-grid prescription is exact and primitive: absolute grid points do not move,
\[
        K_i^{1,\LV}(t)=K_i^0(t).
\]
If the implementation stores local vol in normalized coordinates, do not keep those normalized nodes fixed. Instead preserve the absolute-strike interpolator; the normalized coordinate of a fixed node changes by $x^1=x^0-h_t$.

The Hagan local-vol expansion gives the cheap implied-smile transform. After $F_T^0\mapsto F_T^1$, the new implied volatility at strike $K$ is approximated by the old implied volatility at the displaced strike
\[
        \boxed{\;\sig_1^{\LV}(T,K)\simeq \sig_0(T,K+\Delta F_T)\;} .
\]
Equivalently,
\[
        \boxed{\;\w_1^{\LV}(T,k)\simeq
        \w_0\bigl(T,\ell_T(k,h_T)\bigr)\;},
        \qquad
        \ell_T(k,h_T):=\log\frac{F_T^1\e^k+\Delta F_T}{F_T^0}
        =\log\{\e^{h_T}(\e^k+1)-1\}.
\]
Near ATM and for small moves,
\[
        \ell_T(k,h_T)=k+(1+\e^{-k})h_T+O(h_T^2),
        \qquad
        \ell_T(0,h_T)=2h_T+O(h_T^2),
\]
so
\[
        \w_1^{\LV}(T,k)\simeq \w_0(T,k+2h_T),
        \qquad
        \Delta\sig_{\ATM}^{\LV}=2\kappa_T h_T+O(h_T^2).
\]
The factor two is the usual local-vol ``double skew'' response: the fixed-strike implied vol moves with the old strike-slope, and the ATM strike itself also moves with the forward.

\section*{4. Skew stickiness ratio (SSR)}
The SSR parameter $R_T$ specifies the spot-driven ATM-volatility change in units of the implied skew:
\[
        \boxed{\;\partial_{\log F}\sig_{\ATM}(T)=R_T\kappa_T\;},
        \qquad
        \Delta\sig_{\ATM}(T)\approx R_T\kappa_T h_T .
\]
The canonical values are
\[
        R_T=0 \quad\hbox{sticky-moneyness},\qquad
        R_T=1 \quad\hbox{sticky-strike},\qquad
        R_T\simeq 2 \quad\hbox{sticky-local-vol near ATM}.
\]
The recommended production transform is total-variance horizontal transport,
\[
        \boxed{\;\widetilde\w_1^R(T,k)=\w_0(T,k+R_T h_T)\;} .
\]
This gives
\[
        \sqrt{\widetilde\w_1^R(T,0)/T}
        =\sig_0(T,R_T h_T)
        =\sig_0(T,0)+R_T\kappa_T h_T+O(h_T^2),
\]
and exactly recovers sticky-moneyness at $R_T=0$ and sticky-strike at $R_T=1$; $R_T=2$ matches the Hagan/local-vol ATM expansion.

For finite moves, re-anchor the ATM total variance if an exact linear SSR target is required. Define
\[
        \sig_{\ATM}^{\star}(T):=\sig_0(T,0)+R_T\kappa_T h_T,
        \qquad
        \w_{\ATM}^{\star}(T):=T\bigl(\sig_{\ATM}^{\star}(T)\bigr)^2,
\]
and set
\[
        \boxed{\;\w_1^R(T,k)=\widetilde\w_1^R(T,k)
        +\w_{\ATM}^{\star}(T)-\widetilde\w_1^R(T,0)\;} .
\]
Finally use $\sig_1^R(T,k)=\sqrt{\max(\w_1^R(T,k),\eps)/T}$. One may instead target $\w_0(T,0)+R_Tg_T h_T$ directly.

\textit{Local-vol grid.} Yes: use a barycentric displacement between fixed absolute strike and ATM-relative strike, preferably in log-strike/geometric form. Put
\[
        \boxed{\;\beta_t:=1-\frac{R_t}{2}\;},
        \qquad
        \boxed{\;L_1^R(t,K)\approx L_0\bigl(t,K\e^{-\beta_t h_t}\bigr)\;} .
\]
Equivalently, for old LV node $K_i^0(t)$,
\[
        \boxed{\;K_i^{1,R}(t)=K_i^0(t)\e^{\beta_t h_t}\;},
\]
or, displaying the barycenter explicitly,
\[
        \log K_i^{1,R}
        =\frac{R_t}{2}\log K_i^0
        +\left(1-\frac{R_t}{2}\right)\log\bigl(K_i^0\e^{h_t}\bigr).
\]
Thus $R=0$ moves the grid fully with moneyness, $R=1$ gives the half-grid sticky-strike proxy, and $R=2$ leaves the LV grid fixed. The Hagan/additive version is $K_i^{1,R}\approx K_i^0+(1-R_t/2)\Delta F_t$. An arithmetic barycenter of strike levels is first-order equivalent for small $h_t$, but the log-barycenter preserves positivity and relative spacing and is more stable for production interpolation.

\section*{Summary}
\[
\begin{array}{c|c|c|c}
\toprule
\text{Rule} & \text{Implied total-variance transform} & R_T & \text{LV grid node }K_i^1(t) \\
\midrule
\text{Sticky-moneyness} & \w_0(T,k) & 0 & K_i^0\e^{h_t} \\
\text{Sticky-strike} & \w_0(T,k+h_T) & 1 & K_i^0\e^{h_t/2}\ \text{(proxy)} \\
\text{Sticky-LV} & \w_0(T,\ell_T(k,h_T))\simeq\w_0(T,k+2h_T) & 2 & K_i^0 \\
\text{SSR} & \w_0(T,k+R_T h_T)\ \text{+ optional ATM anchor} & R_T & K_i^0\e^{(1-R_t/2)h_t} \\
\bottomrule
\end{array}
\]
These transforms avoid recalibration for small spot moves and vanilla risk refreshes. For large moves, barrier/exotic risk, or surfaces close to arbitrage bounds, check calendar and butterfly arbitrage and repair or regenerate from the full pricing model if needed.

\vspace{1pt}
{\footnotesize\emph{References.} P. Hagan, D. Kumar, A. Lesniewski, and D. Woodward, ``Managing Smile Risk,'' Wilmott Magazine, 2002. L. Bergomi, ``Smile Dynamics IV,'' 2009, and \emph{Stochastic Volatility Modeling}, Chapman and Hall/CRC, 2015.}
\end{document}
